
%%%%%%%%%%%%%%%%%%%%%%%%%%%%%%%%%%%%%%%%%%%%%%%%%%%%%%%
%%%%%%%%%%%%%%%%%%%%%%%%%%%%%%%%%%%%%%%%%%%%%%%%%%%%%%%
%%%%%%%%%%%%%%%%%%%%%%%%%%%%%%%%%%%%%%%%%%%%%%%%%%%%%%%

\begin{frame}[plain]
  \vfill
  \begin{block}{}
    \centering
    \Huge Bayesian local regression specification checks
  \end{block}
  \vfill
\end{frame}


%%%%%%%%%%%%%%%%%%%%%%%%%%%%%%%%%%%%%%%%%%%%%%%%%%%%%%%
%%%%%%%%%%%%%%%%%%%%%%%%%%%%%%%%%%%%%%%%%%%%%%%%%%%%%%%
%%%%%%%%%%%%%%%%%%%%%%%%%%%%%%%%%%%%%%%%%%%%%%%%%%%%%%%

\begin{frame}[t]{Local specification checks}

% \textbf{Setup}
% % \begin{itemize}
% %  \item Covariates $\x$ (e.g.~race, gender, zip code, age, education level)
% %  \item Responses $\y$ (e.g.~A binary response to ``do you support candidate Z'')
% % \end{itemize}
% %
% \textbf{Survey population:}    Observe \surcol{$(\x_i, y_i)$ for $i = 1, \ldots, \nsur$}\\
% \textbf{Target population:}    Observe \tarcol{$\x_j$ for $j = 1, \ldots, \ntar$}\\[1em]


% \quad$\bullet$ Set $\expect{}{\y \vert \x, \theta} = m(\theta^\trans \x)$ and prior $\p(\theta)$
% (e.g.~Hierarchical logistic regression)\\[1em]

% \quad$\bullet$ Estimate the posterior $\postsur$ with Markov Chain Monte Carlo (MCMC)\\[1em]

% \quad$\bullet$ Take
$\muhatmrp = \tarcol{\meantar \yhat_j}$
where
$\tarcol{\yhat_j} =
    \underbrace{
        \expect{\postsur}{\m(\theta^\trans \tarcol{\x_j})}            
    }_{\textrm{Want to check without re--running MCMC}}
$.


% Standard local approaches:
% %
% \begin{itemize}
% \item Specify an alternative model + test with Bayes factor (CITE LOCAL BF)
% %
% \begin{itemize}
%     \item The model is definitely misspecified, we want to check whether it's \emph{importantly} misspecified
%     \item Bayes factors use prior information (Lindley's paradox)
%     \item Articulating a family of models is inconvenient
%     \item Doesn't assess shrinkage
% \end{itemize}
% %
% \item Use out-of-sample checking (CITE GELMAN)
% \begin{itemize}
%     \item The task is not to predict on a new survey datapoint
% \end{itemize}
% \end{itemize}
%

This model of \emph{survey data} is definitely misspecified.\\[1em]

Why do we care about correct specification? We want to use $\psur(\x)$ to learn about $\ptar(\x)$.\\[1em]

\begin{center}
Correct specification $\Rightarrow$ Regression estimates are invariant to $\psur(\x)$
\end{center}

\textbf{Idea:}  Check invariance to $\psur(\x)$ directly.\\[1em]

Let $\p(\x,\y)$ denote a generic joint density of $\x, \y$.\\
Let $\psurhat$ the empirical distribution on the survey data.
%
\begin{enumerate}
\item Write our estimand as a functional $\mu(\p)$ that takes in a density and returns a MrP estimate
\item Define a covariate density perturbation $\w(\x)$
\item Define $\p_\epsilon(\x, \y) = 
    \p(\x, \y) + \epsilon \w(\x) \p(\y | \x) = 
    \left(\p(\x) + \epsilon \w(\x) \right) \p(\y | \x)$
\item Compute the derivative $\Delta(\w, \p) = \partial \mu(\p_\epsilon) / \partial \epsilon$
\item Plug the empirical distribution into $\Delta(\w, \psurhat)$ and check whether $\approx 0$
\end{enumerate}
%


\end{frame}



%%%%%%%%%%%%%%%%%%%%%%%%%%%%%%%%%%%%%%%%%%%%%%%%%%%%%%%%%%%%%%%%%%%%%%%%%%%
%%%%%%%%%%%%%%%%%%%%%%%%%%%%%%%%%%%%%%%%%%%%%%%%%%%%%%%%%%%%%%%%%%%%%%%%%%%
%%%%%%%%%%%%%%%%%%%%%%%%%%%%%%%%%%%%%%%%%%%%%%%%%%%%%%%%%%%%%%%%%%%%%%%%%%%

\def\sensitivitysetup{
Consider a family of distributions
$
% $$
\peps(\theta) 
%:= \frac{\p(\theta) \exp(\epsilon \h(\theta))}{\int \p(\theta') \exp(\epsilon \h(\theta')) d\theta'}
\propto \p(\theta) \exp(\epsilon \h(\theta))
% $$
$
and a test function $\phi(\theta)$.
%  where we are interested in 
% $$
% \expect{\peps(\theta)}{\phi(\theta)}.
% $$
Then, by interchange of integration and differentiation,
$$
\fracat{\partial\expect{\peps(\theta)}{\phi(\theta)}}
       {\partial \epsilon}{\epsilon = 0} =
\cov{\p(\theta)}{\phi(\theta), \h(\theta)}.
$$
}


\begin{frame}[t]{Bayesian sensitivity analysis}

\sensitivitysetup{}

\hrulefill

Let the regression log likelihood be $\ell(\y \vert \x, \theta)$, so that
$$
\post \propto \exp\left(\sumsur \ell(\y_i \vert \x_i, \theta) \right)
\textrm{ and }
\posteps \propto 
%\exp\left(\sumsur (1 + \epsilon \w(\x_i))\ell(\y_i \vert \x_i, \theta) \right)
\post \exp\left(\sumsur \epsilon \w(\x_i) \ell(\y_i \vert \x_i, \theta) \right)
$$
%
\begin{itemize}
\item $\p(\theta) = \post$
\item $\phi(\theta) = \meantar \m(\theta^\trans \x_j)$  (so that $\muhatmrp = \expect{\post}{\phi(\theta)})$
\item $\h(\theta) = \sumsur \w(\x_i) \ell(\y_i \vert \x_i, \theta)$ \Rightarrow
\end{itemize}
%
Our diagnostic is then given by
$$
\begin{aligned}
    \Delta(\w, \psurhat) :=
    \fracat{\partial \muhatmrp}{\partial \epsilon}{\epsilon = 0} =
    \cov{\post}{\meantar \m(\theta^\trans \x_j), \sumsur \w(\x_i) \ell(\y_i \vert \x_i, \theta)}
\end{aligned}
$$
We expect $\Delta(\w, \psurhat)$ to be close to zero when the model is correctly
specified and posterior estimation is accurate.


\end{frame}







% %%%%%%%%%%%%%%%%%%%%%%%%%%%%%%%%%%%%%%%%%%%%%%%%%%%%%%%
% %%%%%%%%%%%%%%%%%%%%%%%%%%%%%%%%%%%%%%%%%%%%%%%%%%%%%%%
% %%%%%%%%%%%%%%%%%%%%%%%%%%%%%%%%%%%%%%%%%%%%%%%%%%%%%%%


\def\generalizedpost{
Define the ``generalized posterior'' functional \citep{giordano:2023:bayesij}:
%
\begin{align*}
    %
    \T(\p, N) := \frac{
        \int \g(\theta) \exp\left( 
            N \int \ell(\yn \vert \x, \theta) \p(\xn, \yn) d\xn d\yn
        \right)
        \prior(\theta) d\theta
    }
    {
    \int \exp\left( 
            N \int \ell(\yn \vert \x, \theta) \p(\xn, \yn) d\xn d\yn
    \right)
    \prior(\theta) d\theta
    }.
    %
\end{align*}
%
}



% %%%%%%%%%%%%%%%%%%%%%%%%%%%%%%%%%%%%%%%%%%%%%%%%%%%%%%%%%%%%%%%%%%%%%%%%%%%
% %%%%%%%%%%%%%%%%%%%%%%%%%%%%%%%%%%%%%%%%%%%%%%%%%%%%%%%%%%%%%%%%%%%%%%%%%%%
% %%%%%%%%%%%%%%%%%%%%%%%%%%%%%%%%%%%%%%%%%%%%%%%%%%%%%%%%%%%%%%%%%%%%%%%%%%%

% \begin{frame}[t]{Bayesian posterior functionals}

% Supppose we are interested in $\expect{\post}{\g(\theta)}$.
% Let the regression log likelihood be $\ell(\y \vert \x, \theta)$.\\[1em]

% \generalizedpost{}

% Let $\psurhat$ denote the empirical distribution over $\x_i, \y_i$.  Then
% %
% \begin{align*}
%     \T(\psurhat, \nsur) =
%     \frac{
%         \int \g(\theta) \exp\left( \nsur \meansur \ell(\y_i \vert \x_i, \theta) \right)
%             \prior(\theta) d\theta
%     }
%     {
%     \int \exp\left( \nsur \meansur \ell(\y_i \vert \x_i, \theta) \right)
%         \prior(\theta) d\theta
%     } = \expect{\post}{\g(\theta)}.
% \end{align*}


% Consider a smooth path of densities $\epsilon \mapsto \p^\epsilon$. 
% The directional derivative in general is
% $$
% \begin{aligned}
% \fracat{\partial \T(\p^\epsilon, N)}{\partial \epsilon}{\epsilon=0} ={}&
%     \cov{\p(\theta \vert \p, N)}{
%         \g(\theta), 
%         N \int \ell(\yn \vert \xn, \theta) \w(d \xn)
%         \fracat{\partial \log \p_\epsilon(\yn \vert \xn)}
%                {\partial \epsilon}{\epsilon = 0} d\xn d\yn
%     }
% \\\Rightarrow
% \fracat{\partial \T({\psurhat}^\epsilon, \nsur)}{\partial \epsilon}{\epsilon=0} ={}&
%     \cov{\post}{
%         \g(\theta), 
%         \sumsur \ell(\y_i \vert \x_i, \theta)
%         \fracat{\partial \log {\psurhat}^\epsilon(\y_i \vert \x_i)}
%                {\partial \epsilon}{\epsilon = 0}
%     }
% \end{aligned}
% $$

% \end{frame}



% %%%%%%%%%%%%%%%%%%%%%%%%%%%%%%%%%%%%%%%%%%%%%%%%%%%%%%%%%%%%%%%%%%%%%%%%%%%
% %%%%%%%%%%%%%%%%%%%%%%%%%%%%%%%%%%%%%%%%%%%%%%%%%%%%%%%%%%%%%%%%%%%%%%%%%%%
% %%%%%%%%%%%%%%%%%%%%%%%%%%%%%%%%%%%%%%%%%%%%%%%%%%%%%%%%%%%%%%%%%%%%%%%%%%%

% \begin{frame}[t]{Bayesian posterior functionals}

% We're interested in $\g(\theta) = \meantar \m(\theta^\trans \x_j)$.

% For some function $\w(\xn)$, set
% $$
% \psur^\epsilon(\xn, yn) = 
%     \psur(\xn, \yn) + \epsilon \w(d \xn) \psur(\xn, \yn) =
%     \left( \psur(\xn) + \epsilon \w(\xn) \right) \psur(\yn \vert \xn).
% $$
% This is a shift in the marginal distribution of $\x$ in the direction $\w(\x)$,
% with\\[1em]
% $$
% \begin{aligned}
%     \fracat{\partial \psur^\epsilon(\yn \vert \xn)}
%                {\partial \epsilon}{\epsilon = 0} =
%                \w(\xn) \psur(\yn \vert \xn).
% \end{aligned}
% $$

% Our diagnostic is then given by
% $$
% \Delta(\w, \psurhat) =
%     \cov{\post}{
%         \meantar \m(\theta^\trans \x_j),
%         \sumsur \w(\x_i) \ell(\y_i \vert \x_i, \theta)
%         }
% $$
% %
% We expect $\Delta(\w, \psurhat)$ to be close to zero when the model is correctly
% specified and posterior estimation is accurate.
% %
% \end{frame}



%%%%%%%%%%%%%%%%%%%%%%%%%%%%%%%%%%%%%%%%%%%%%%%%%%%%%%%
%%%%%%%%%%%%%%%%%%%%%%%%%%%%%%%%%%%%%%%%%%%%%%%%%%%%%%%
%%%%%%%%%%%%%%%%%%%%%%%%%%%%%%%%%%%%%%%%%%%%%%%%%%%%%%%


\begin{frame}{Simulation}
    \begin{minipage}{0.95\textwidth}
    % \only<1>{\ShiftGraphZero{}}
    % \only<2>{\ShiftGraphOne{}}
    \only<1>{\OttoGraphOne{}}
    \end{minipage}

Simulation data: Logistic regression with five active regressors, each
with covariate shift.
\onslide<1->{
$$
\begin{aligned}
    \textrm{Model 0} \textrm{ (correct)}:{}&  \texttt{\y} \sim \textrm{Logistic}(\texttt{1 + X1 + X2 + X3 + X4 + X5})\\
    \textrm{Model 1} \textrm{ (misspecified)}:{}&  \texttt{\y} \sim \textrm{Logistic}(\texttt{1 + X3 + X4 + X5})\\
\end{aligned}
$$
}
\onslide<1->{
Compute diagnostic with $\w(\x) = \ind{\texttt{X2} > 0.878}$
}

\end{frame}




